% Rachel: demographic parity under group-dependent acceptance.
% \input from a1.tex after sections/setup (shared model, E_c, r_c, ceiling model, Claim 1). Derivations: docs/parity.md.

\textbf{Parity objectives.} Following Fang et al., the \emph{DP gap} is $\Delta_{DP}=\max_c r_c-\min_c r_c$; we also use \emph{$\varepsilon$-DP} ($\Delta_{DP}\le\varepsilon$), the ratio form $\min_c r_c/\max_c r_c\ge\tau$, and the scalarized objective $F_\lambda=B/E-\lambda\Delta_{DP}$. $F_\lambda$ is neither monotone nor submodular: a seed that raises one group's rate can widen the gap by more than it adds.

\emph{Claim 2 (leveling down).} Past $k^*$, $L$ is saturated ($r_L=a_L$), so parity caps $r_H\le a_L$. DP's total stays at $a_L N$, while lexicographic maximin (FWRRS's tie-break) sends extra seeds to $H$ and gains $\min\bigl(b\,a_H(k-k^*),\,(a_H-a_L)E_H\bigr)$ more. DP's PoF keeps growing with $k$ up to $r\,w_H+w_L$, which is unbounded only as $a_L\to0$. Plain maximin is indifferent past $k^*$.

\emph{Claim 3 (relaxations).} Past $k^*$, $\varepsilon$-DP lets $r_H$ rise to $\min(a_L+\varepsilon,a_H)$ and the ratio form to $\min(a_L/\tau,a_H)$, so total $\le a_L E_L+\min(\cdot)\,E_H$.

Figs.~\ref{fig:parity-accept}--\ref{fig:parity-pareto} test Claims 1--3.
